The previous section compares cards against a single pooled vector, which leaves
an obvious objection open: cards store more vectors per document, so perhaps
capacity alone explains the gain and any decomposition would do. This section
answers that objection with a control.

\paragraph{Setting.}
LIMIT~\cite{weller2025limit} is constructed to expose the representational
ceiling of single-vector retrieval. Each document names a person and lists the
things they like; each query asks who likes one specific thing. The small variant
has 46 documents, 44.0 attributes each, and 1{,}000 queries, so the task looks
trivial and yet defeats dense retrievers. It also hands us an unambiguous aspect
decomposition, which makes it the sharpest available test of whether alignment
matters. We compare five systems, all scored at document level: BM25; a pooled
dense vector; the same text cut into 30-word windows with 10-word overlap
(186 chunks); the same text cut into per-attribute cards (2{,}024 cards); and
reciprocal rank fusion of the card system with BM25.

\begin{table}[t]
\caption{LIMIT-small. Deltas are against the pooled dense baseline, paired
permutation with Holm correction over 1{,}000 queries. Chunking and carding both
raise the embedding count from one per document to many; only carding recovers
the lost performance.}
\label{tab:limit}
\centering
\begin{tabular}{lrrrrl}
\toprule
system & nDCG@10 & R@2 & R@10 & $\Delta$ vs pooled & $p$ \\
\midrule
BM25                       & 0.997 & 0.991 & 1.000 & $+0.683$ & $<0.001$ \\
dense, pooled              & 0.314 & 0.162 & 0.480 & baseline & \\
dense, 30-word chunks      & 0.390 & 0.225 & 0.549 & $+0.076$ & $<0.001$ \\
dense, per-aspect cards    & 0.988 & 0.965 & 0.997 & $+0.674$ & $<0.001$ \\
cards $\oplus$ BM25 (RRF)  & 0.997 & 0.990 & 1.000 & $+0.683$ & $<0.001$ \\
\bottomrule
\end{tabular}
\end{table}

\paragraph{Alignment, not count.}
Table~\ref{tab:limit} separates the two explanations. Chunking multiplies the
embedding budget and recovers 0.076 nDCG@10. Carding multiplies the same budget
and recovers 0.674, roughly nine times as much, from identical source text under
an identical encoder. The difference between the two is only where the boundaries
fall: a 30-word window spans several attributes and dilutes each of them in the
same way the pooled vector does, whereas a card corresponds to exactly the unit a
query names. Capacity is necessary and not sufficient; the decomposition has to
match the query distribution.

\paragraph{What this result does not show.}
LIMIT supplies a perfect decomposition for free, and its queries are effectively
exact-attribute lookups, which is why BM25 wins outright and the hybrid only
matches it. Real corpora arrive undecomposed, the taxonomy must be designed and
may be misaligned, and the honest reading of this table is as an upper bound on
what alignment can buy rather than as an estimate of what it will buy. We also
note the cost: 44 cards per document here, far above the small constant that the
pattern implies elsewhere, so the storage multiplier is a property of the corpus
and must be reported per corpus rather than assumed.
